\begin{lemma}[The weak gradient of a compactly supported function has mean zero]
Let $n\ge1$ and $q\ge1$, let $x_0\in\R^n$ and $R>0$, and let
$\varphi\in W^{1,q}(B_R(x_0))$ with weak gradient $g_\varphi$
(\noderef{MemW1pOn}) be such that $\varphi$ has compact support and
$\operatorname{tsupport}\varphi\subseteq B_R(x_0)$.  Then for every
$i\in\{1,\dots,n\}$ the component $g_\varphi^{\,i}$ is integrable on $B_R(x_0)$
and
\[
\int_{B_R(x_0)}g_\varphi^{\,i}\,dx\ =\ 0 .
\]

This is the weak-gradient analogue of ``the integral of a derivative of a
compactly supported function vanishes''.  Both conclusions are needed together:
the vanishing of the integral is empty information without the integrability
that makes the integral the integral of $g_\varphi^{\,i}$ rather than the
Bochner convention's $0$.  The hypotheses are exactly the ones carried by the
test pairs of \noderef{IsWeakSolutionOn} at a ball, which is where the statement
is consumed: a weak-solution identity is unchanged when a constant vector is
subtracted from the field precisely because of this lemma.  Nothing here
mentions an operator, a solution or a structure condition --- it is a statement
about the weak-gradient predicate alone.

Note that the support condition is on the \emph{closed} support
$\operatorname{tsupport}\varphi=\overline{\{\varphi\ne0\}}$, and that $\varphi$
is a function on all of $\R^n$ whose closed support happens to lie in the ball;
this is the same convention as \noderef{MemW1pOn}'s consumers use.  The
hypothesis $q\ge1$ is used only to get $L^{q}(B_R)\subseteq L^{1}(B_R)$ on the
ball, which has finite measure; no upper bound on $q$ is needed.
\end{lemma}
\begin{proof}
Write $B=B_R(x_0)$ and fix $i$.

\emph{Step 1: integrability.}  By \noderef{MemW1pOn}, $g_\varphi\in L^{q}(B)$.
The ball $B$ has finite Lebesgue measure, so $L^{q}(B)\subseteq L^{1}(B)$
(Preamble: H\"older on a set of finite measure, the stated inclusion
$L^{q}(E)\subseteq L^{1}(E)$ for $\mu(E)<\infty$ and $q\ge1$), and hence
$g_\varphi$ is integrable on $B$.  For every $v\in\R^n$ and every $i$ one has
$|v^i|\le\|v\|$, because $(v^i)^2\le\sum_m(v^m)^2=\|v\|^2$ and both
$|v^i|$ and $\|v\|$ are nonnegative.  So
$|g_\varphi^{\,i}|\le\|g_\varphi\|$ pointwise, and the measurable function
$g_\varphi^{\,i}$ --- measurable because it is the composition of $g_\varphi$
with the continuous linear coordinate projection --- is dominated by an
integrable function on $B$; hence $g_\varphi^{\,i}$ is integrable on $B$
(Preamble: domination).

\emph{Step 2: extending the pair to $\R^n$.}  Apply
\noderef{WeakGradientZeroExtend} with $n$, the exponent $q$, the open set
$\Omega=\R^n$, the centre $x_0$, the radius $R$ and the pair
$(\varphi,g_\varphi)$.  Its hypotheses are $1\le n$, $1\le q$, the openness of
$\R^n$, the inclusion $B\subseteq\R^n$, the membership
$\varphi\in W^{1,q}(B)$ with weak gradient $g_\varphi$, the compactness of the
support of $\varphi$ and $\operatorname{tsupport}\varphi\subseteq B$; all are in
hand.  Its conclusion is that
$\varphi\in W^{1,q}(\R^n)$ with weak gradient $\mathbf1_Bg_\varphi$, the field
equal to $g_\varphi$ on $B$ and to $0$ elsewhere (\noderef{MemW1pOn}).  In
particular $\mathbf1_Bg_\varphi$ is a weak gradient of $\varphi$ on $\R^n$
(\noderef{HasWeakGradientOn}).

\emph{Step 3: a cut-off that is $1$ on a ball containing $B$.}  Put
$M=\|x_0\|+R+1>0$.  Then $B\subseteq B_M(0)$: if $|y-x_0|<R$ then
$|y|\le|y-x_0|+|x_0|<R+\|x_0\|<M$.  Apply \noderef{StandardRadialCutoff} with
$n$ --- its hypothesis is $1\le n$ --- at the centre $0$ and the radii
$\rho_0=M$ and $\rho_1=M+1$, whose hypotheses $0<\rho_0$ and $\rho_0<\rho_1$
hold.  Let $\zeta$ be the resulting function: $\zeta$ is smooth, has compact
support, $\operatorname{tsupport}\zeta\subseteq B_{M+1}(0)$, and $\zeta=1$ at
every point of $\overline B_M(0)$.  (The gradient bound that node also supplies
is not used here.)

\emph{Step 4: the integration-by-parts identity at $\zeta$.}  The function
$\zeta$ is an admissible test function for \noderef{HasWeakGradientOn} on
$\R^n$: it is smooth, compactly supported, and
$\operatorname{tsupport}\zeta\subseteq\R^n$ trivially.  So Step 2 gives
\[
\int_{\R^n}\varphi\,D_i\zeta\,dx\ =\ -\int_{\R^n}(\mathbf1_Bg_\varphi)^i\,\zeta\,dx .
\]

\emph{Step 5: the left-hand side is $0$.}  We show the integrand vanishes at
every point $x$.  If $x\notin\operatorname{tsupport}\varphi$ then
$\varphi(x)=0$ and the product is $0$.  If
$x\in\operatorname{tsupport}\varphi$ then $x\in B\subseteq B_M(0)$, and $B_M(0)$
is an open set on which $\zeta$ is identically $1$ --- by Step 3, since
$B_M(0)\subseteq\overline B_M(0)$.  A function that is constant on a
neighbourhood of $x$ has vanishing derivative at $x$, so $D_i\zeta(x)=0$ and the
product is $0$ again.  Hence the left-hand side of Step 4 is the integral of the
zero function, which is $0$.

\emph{Step 6: the right-hand side is the integral we want.}  We show that at
every point $x$,
\[
(\mathbf1_Bg_\varphi(x))^i\,\zeta(x)
=\bigl(\mathbf1_B\,g_\varphi^{\,i}\bigr)(x) ,
\]
the right-hand side meaning $g_\varphi^{\,i}(x)$ for $x\in B$ and $0$ otherwise.
If $x\in B$ then $\mathbf1_Bg_\varphi(x)=g_\varphi(x)$ and, since
$B\subseteq B_M(0)\subseteq\overline B_M(0)$, also $\zeta(x)=1$ by Step 3, so
both sides are $g_\varphi^{\,i}(x)$.  If $x\notin B$ then
$\mathbf1_Bg_\varphi(x)=0$, so the left-hand side is $0\cdot\zeta(x)=0$, and the
right-hand side is $0$ as well.  Therefore
\[
\int_{\R^n}(\mathbf1_Bg_\varphi)^i\,\zeta\,dx
=\int_{\R^n}\mathbf1_B\,g_\varphi^{\,i}\,dx
=\int_Bg_\varphi^{\,i}\,dx ,
\]
the last equality because $B$ is measurable and the integral over $\R^n$ of the
zero-extension of a function off a measurable set is its integral over that set.

\emph{Conclusion.}  Combining Steps 4, 5 and 6,
$0=-\int_Bg_\varphi^{\,i}\,dx$, i.e.\ $\int_Bg_\varphi^{\,i}\,dx=0$.  With Step
1 this is the assertion of the lemma.
\end{proof}
